\chapter{The Kitaev Honeycomb Model}
\label{chap:kitaev-honeycomb-model}

The toric code and the quantum double models are exactly solvable because their
Hamiltonians are sums of commuting projectors.  The Kitaev honeycomb model is
solvable by a different, but closely related, mechanism.  It is not a
commuting-projector Hamiltonian.  Instead, its spin degrees of freedom can be
represented by Majorana fermions coupled to a static \(\ZZ_2\) gauge field.  In
each fixed gauge-flux sector, the Hamiltonian reduces to a free Majorana hopping
problem.

This chapter therefore serves as the bridge between two ideas developed earlier
in these notes:
\begin{enumerate}[label=(\roman*)]
    \item free-fermion/BdG solvability, as in the transverse-field XY chain and
    the Kitaev chain;
    \item gauge-sector solvability, as in the toric code and finite-group
    quantum double models.
\end{enumerate}
The special feature of the honeycomb model is that these two structures coexist
in a two-dimensional spin system.  The resulting phase diagram contains gapped
Abelian topological phases, a gapless Majorana spin liquid, and, after a
small time-reversal-breaking perturbation, a gapped non-Abelian Ising-anyon
phase.

\section{Honeycomb lattice and bond-dependent exchange}
\label{sec:kitaev-honeycomb-hamiltonian}

Consider a honeycomb lattice.  Each vertex carries a spin-\(1/2\) degree of
freedom, and the nearest-neighbour bonds are colored by three labels
\(\alpha=x,y,z\).  Every site touches exactly one bond of each type.  The Kitaev
honeycomb Hamiltonian is
\begin{equation}
\label{eq:kitaev-honeycomb-spin-hamiltonian}
    H_{\rm KHC}
    =
    -J_x\sum_{\langle ij\rangle_x}\sigma_i^x\sigma_j^x
    -J_y\sum_{\langle ij\rangle_y}\sigma_i^y\sigma_j^y
    -J_z\sum_{\langle ij\rangle_z}\sigma_i^z\sigma_j^z .
\end{equation}
Here \(\langle ij\rangle_\alpha\) denotes a nearest-neighbour bond of type
\(\alpha\).  Unless otherwise stated, we take
\begin{equation}
\label{eq:kitaev-couplings-positive}
    J_x,J_y,J_z\geq0 .
\end{equation}
Changing signs of the couplings can often be absorbed by sublattice-dependent
spin rotations, but keeping the positive-coupling convention makes the phase
diagram transparent.

\begin{remark}[Why the model is not trivially solvable]
Each local term in \cref{eq:kitaev-honeycomb-spin-hamiltonian} is an Ising
exchange, but different bonds use different spin components.  Therefore terms
sharing one site usually do not commute.  For example, if an \(x\)-bond and a
\(y\)-bond meet at site \(i\), then the local Pauli matrices
\(\sigma_i^x\) and \(\sigma_i^y\) anticommute.  Thus the model is not a
commuting-projector model.  Its exact solvability is more subtle: noncommuting
spin interactions become quadratic Majorana hopping terms in a background of
static \(\ZZ_2\) link variables.
\end{remark}

The honeycomb lattice is bipartite.  We denote the two sublattices by
\(A\) and \(B\).  Choose every bond orientation from \(A\) to \(B\).  A unit cell
contains one \(A\)-site and one \(B\)-site.  Let \(\bm a_1\) and \(\bm a_2\) be
primitive lattice vectors chosen so that, in the vortex-free gauge used below,
the three bonds from an \(A\)-site at cell \(\bm r\) connect to
\begin{equation}
\label{eq:honeycomb-neighbour-convention}
    B_{\bm r},
    \qquad
    B_{\bm r+\bm a_1},
    \qquad
    B_{\bm r+\bm a_2},
\end{equation}
with couplings \(J_z\), \(J_x\), and \(J_y\), respectively.  This choice fixes
the later Bloch function
\begin{equation}
\label{eq:kitaev-bloch-function-def-preview}
    f(\bm k)=J_z+J_x\ee^{\ii \bm k\cdot \bm a_1}
             +J_y\ee^{\ii \bm k\cdot \bm a_2} .
\end{equation}
Other conventions only multiply \(f(\bm k)\) by phase factors or relabel the
momenta; the spectrum is unchanged.

\section{Plaquette fluxes in the spin language}
\label{sec:kitaev-spin-fluxes}

Before introducing Majorana fermions, one can already identify conserved
plaquette operators directly in the spin model.  Let \(p\) be a hexagonal
plaquette and label its six vertices cyclically by \(1,2,3,4,5,6\).  With the
usual bond coloring around a hexagon, define
\begin{equation}
\label{eq:kitaev-plaquette-spin-operator}
    W_p
    =
    \sigma_1^x\sigma_2^y\sigma_3^z
    \sigma_4^x\sigma_5^y\sigma_6^z,
\end{equation}
up to a cyclic relabeling determined by the local bond convention.  The precise
sequence of \(x,y,z,x,y,z\) may be shifted, but the operator is always the
product of the spin component not associated with the outgoing external bond at
each vertex of the plaquette.

These operators obey
\begin{equation}
\label{eq:kitaev-flux-algebra-spin}
    W_p^2=\id,
    \qquad
    [W_p,W_{p'}]=0,
    \qquad
    [W_p,H_{\rm KHC}]=0 .
\end{equation}
Thus every plaquette carries a conserved \(\ZZ_2\) flux eigenvalue
\begin{equation}
\label{eq:kitaev-flux-eigenvalue}
    W_p=w_p,
    \qquad
    w_p=\pm1 .
\end{equation}
A plaquette with \(w_p=-1\) is called a vortex or \(\ZZ_2\) flux excitation.

\begin{remark}[Fluxes versus commuting projectors]
The fluxes \(W_p\) commute with the Hamiltonian, but the Hamiltonian itself is
not a sum of the \(W_p\).  They are conserved labels of superselection sectors,
not the complete set of local energy projectors.  After fixing the flux sector,
one still has to diagonalize itinerant Majorana fermions.
\end{remark}

On a closed surface, the plaquette fluxes are not all independent.  For example,
on a torus,
\begin{equation}
\label{eq:kitaev-product-flux-constraint}
    \prod_p W_p=\id,
\end{equation}
so vortices are created in pairs on a closed lattice.  There are also nonlocal
Wilson-loop sectors associated with the two noncontractible cycles of the torus.
These global sectors will later become important for the finite-size ground
state degeneracy.

\section{Four-Majorana representation}
\label{sec:kitaev-four-majorana-representation}

The exact solution begins by enlarging the Hilbert space.  At every site \(i\),
introduce four Hermitian Majorana operators
\begin{equation}
\label{eq:kitaev-four-majoranas}
    b_i^x,
    \qquad
    b_i^y,
    \qquad
    b_i^z,
    \qquad
    c_i,
\end{equation}
satisfying
\begin{equation}
\label{eq:kitaev-majorana-algebra}
    \{b_i^\alpha,b_j^\beta\}=2\delta_{ij}\delta_{\alpha\beta},
    \qquad
    \{c_i,c_j\}=2\delta_{ij},
    \qquad
    \{b_i^\alpha,c_j\}=0 .
\end{equation}
The spin operators are represented as
\begin{equation}
\label{eq:kitaev-spin-majorana-representation}
    \sigma_i^\alpha=\ii b_i^\alpha c_i,
    \qquad
    \alpha=x,y,z .
\end{equation}

This representation is redundant.  Four Majoranas span a four-dimensional local
Hilbert space, whereas one spin-\(1/2\) has dimension two.  The physical
subspace is selected by the local gauge constraint
\begin{equation}
\label{eq:kitaev-local-gauge-constraint}
    D_i\ket{\psi}_{\rm phys}=\ket{\psi}_{\rm phys},
    \qquad
    D_i=b_i^x b_i^y b_i^z c_i .
\end{equation}
The operator \(D_i\) satisfies \(D_i^2=\id\).  It generates a local
\(\ZZ_2\) gauge transformation:
\begin{equation}
\label{eq:kitaev-gauge-flip-majoranas}
    D_i b_i^\alpha D_i=-b_i^\alpha,
    \qquad
    D_i c_i D_i=-c_i,
\end{equation}
while leaving Majoranas on other sites invariant.  Since each spin operator
\(\sigma_i^\alpha=\ii b_i^\alpha c_i\) contains two Majoranas at site \(i\), it
is gauge invariant.

The projector onto the physical spin Hilbert space is
\begin{equation}
\label{eq:kitaev-physical-projector}
    P_{\rm phys}
    =
    \prod_i \frac{1+D_i}{2} .
\end{equation}
On finite lattices, especially with periodic boundary conditions, this projection
must be imposed carefully because it constrains the allowed fermion parity in
each gauge sector.  For deriving the bulk spectrum and the phase diagram, it is
usually sufficient to diagonalize in the enlarged Majorana space and then project
physical states at the end.

\begin{proposition}[Pauli algebra on the physical subspace]
The operators \(\sigma_i^\alpha=\ii b_i^\alpha c_i\) satisfy the spin-\(1/2\)
Pauli algebra on the subspace \(D_i=+1\):
\begin{equation}
\label{eq:kitaev-pauli-algebra-physical}
    \sigma_i^\alpha\sigma_i^\beta
    =
    \delta_{\alpha\beta}\id
    +\ii\sum_{\gamma}\epsilon_{\alpha\beta\gamma}\sigma_i^\gamma .
\end{equation}
\end{proposition}

\begin{proof}
For \(\alpha=\beta\), one has
\((\ii b_i^\alpha c_i)^2=\id\).  For distinct \(\alpha,\beta,\gamma\),
\begin{equation}
\label{eq:kitaev-pauli-proof-product}
    \sigma_i^\alpha\sigma_i^\beta
    =
    (\ii b_i^\alpha c_i)(\ii b_i^\beta c_i)
    =
    b_i^\alpha b_i^\beta .
\end{equation}
Using \(D_i=b_i^x b_i^y b_i^z c_i=+1\), one obtains
\begin{equation}
\label{eq:kitaev-pauli-proof-constraint-use}
    b_i^\alpha b_i^\beta
    =
    \ii\epsilon_{\alpha\beta\gamma}\,\ii b_i^\gamma c_i
    =
    \ii\epsilon_{\alpha\beta\gamma}\sigma_i^\gamma,
\end{equation}
with the sign fixed by the ordering \((x,y,z)\).  This gives
\cref{eq:kitaev-pauli-algebra-physical}.
\end{proof}

\section{Static \texorpdfstring{\(\ZZ_2\)}{Z2} gauge field}
\label{sec:kitaev-static-z2-gauge-field}

Substituting \cref{eq:kitaev-spin-majorana-representation} into the spin
Hamiltonian gives
\begin{align}
\label{eq:kitaev-majorana-hamiltonian-derivation}
    -J_\alpha\sigma_i^\alpha\sigma_j^\alpha
    &=-J_\alpha(\ii b_i^\alpha c_i)(\ii b_j^\alpha c_j)  \notag\\
    &=\ii J_\alpha u_{ij} c_i c_j,
    \qquad
    \langle ij\rangle_\alpha,
\end{align}
where the link operator is
\begin{equation}
\label{eq:kitaev-link-variable}
    u_{ij}=\ii b_i^\alpha b_j^\alpha,
    \qquad
    \langle ij\rangle_\alpha .
\end{equation}
The orientation convention is important: throughout this chapter \(i\in A\) and
\(j\in B\), so \(u_{ji}=-u_{ij}\).  Since \(u_{ij}^2=\id\), its eigenvalues are
\begin{equation}
\label{eq:kitaev-link-eigenvalues}
    u_{ij}=\pm1 .
\end{equation}
The Hamiltonian becomes
\begin{equation}
\label{eq:kitaev-majorana-hamiltonian-static-gauge}
    H_{\rm KHC}
    =
    \ii\sum_{\langle ij\rangle_\alpha}J_\alpha u_{ij}c_i c_j .
\end{equation}

\begin{proposition}[Static link variables]
The link operators satisfy
\begin{equation}
\label{eq:kitaev-u-commutation}
    [u_{ij},u_{kl}]=0,
    \qquad
    [u_{ij},H_{\rm KHC}]=0 .
\end{equation}
Therefore each \(u_{ij}\) is a conserved \(\ZZ_2\) variable in the enlarged
Majorana Hilbert space.
\end{proposition}

\begin{proof}
Two link variables either act on disjoint Majoranas or share one site but use
Majoranas of different bond types.  In both cases the total number of Majorana
interchanges needed to swap the two bilinears is even, so they commute.

To prove \([u_{ij},H_{\rm KHC}]=0\), note first that \(u_{ij}\) trivially
commutes with the corresponding bond term
\(\ii J_\alpha u_{ij}c_i c_j\).  For any other bond term, either the two bonds
are disjoint, or they share one site but involve different \(b\)-Majoranas and a
single \(c\)-Majorana at that site.  Again the relevant operator products differ
by an even number of Majorana interchanges.  Hence \(u_{ij}\) commutes with
every term in \cref{eq:kitaev-majorana-hamiltonian-static-gauge}.
\end{proof}

The conserved \(u_{ij}\) are not themselves gauge-invariant observables.  Under
the local gauge transformation generated by \(D_i\), all link variables adjacent
to site \(i\) change sign:
\begin{equation}
\label{eq:kitaev-link-gauge-transformation}
    u_{ij}\longmapsto -u_{ij},
    \qquad
    u_{ki}\longmapsto -u_{ki} .
\end{equation}
Gauge-invariant information is contained in products around closed loops.  For a
plaquette \(p\), define
\begin{equation}
\label{eq:kitaev-majorana-flux}
    W_p
    =
    \prod_{\langle ij\rangle\in\partial p} u_{ij},
\end{equation}
where the product follows the oriented boundary, with the sign convention fixed
by the chosen bond orientations.  This operator is the Majorana representation
of the spin flux in \cref{eq:kitaev-plaquette-spin-operator}.  The eigenvalue
\(W_p=+1\) is called the vortex-free or zero-flux sector, and \(W_p=-1\) is a
\(\ZZ_2\) vortex.

\section{Reduction to free Majorana fermions}
\label{sec:kitaev-free-majorana-reduction}

Because all \(u_{ij}\) commute with the Hamiltonian, the Hilbert space splits
into static gauge sectors.  In a fixed configuration
\(u=\{u_{ij}=\pm1\}\), the Hamiltonian is quadratic:
\begin{equation}
\label{eq:kitaev-fixed-sector-hamiltonian}
    H[u]
    =
    \ii\sum_{\langle ij\rangle_\alpha}J_\alpha u_{ij}c_i c_j .
\end{equation}
Equivalently, it can be written in the canonical Majorana quadratic form
\begin{equation}
\label{eq:kitaev-majorana-quadratic-form}
    H[u]
    =
    \frac{\ii}{4}\sum_{i,j}A_{ij}[u]c_i c_j,
    \qquad
    A_{ij}[u]=-A_{ji}[u]\in\RR .
\end{equation}
For an \(\alpha\)-bond oriented from \(i\in A\) to \(j\in B\), the matrix entries
are
\begin{equation}
\label{eq:kitaev-a-matrix-entry}
    A_{ij}[u]=2J_\alpha u_{ij},
    \qquad
    A_{ji}[u]=-2J_\alpha u_{ij} .
\end{equation}
Diagonalizing the real antisymmetric matrix \(A[u]\) gives pairs of eigenvalues
\(\pm\epsilon_n\) of \(\ii A[u]\), with \(\epsilon_n\geq0\).  The Hamiltonian in
that sector has the form
\begin{equation}
\label{eq:kitaev-diagonal-majorana-spectrum}
    H[u]
    =
    \sum_{n}\epsilon_n\left(f_n^\dagger f_n-\frac{1}{2}\right),
\end{equation}
up to the physical projection discussed in \cref{eq:kitaev-physical-projector}.
Thus the original interacting spin model is solved by the following sequence:
\begin{align}
\label{eq:kitaev-solution-sequence}
    \text{spin model}
    &\longrightarrow
    \text{Majoranas plus static }\ZZ_2\text{ gauge field}
    \notag\\
    &\longrightarrow
    \text{free Majorana band problem in each flux sector}
    \notag\\
    &\longrightarrow
    \text{physical projection} .
\end{align}

\begin{principle}[Majorana--gauge solvability]
A spin model is Majorana--gauge solvable if, after an enlarged Majorana
representation and a local gauge constraint, its Hamiltonian becomes quadratic in
itinerant Majoranas with coefficients given by conserved \(\ZZ_2\) link
variables.  The Kitaev honeycomb model is the canonical two-dimensional example
of this mechanism.
\end{principle}

\section{Vortex-free sector and Bloch Hamiltonian}
\label{sec:kitaev-vortex-free-bloch}

For the honeycomb lattice with positive couplings, the ground state lies in the
vortex-free sector.  This is a consequence of a flux theorem for reflection-
symmetric planar lattices.  We will not prove the theorem here; instead, we use
it to obtain the bulk phase diagram.

Choose the translationally invariant gauge
\begin{equation}
\label{eq:kitaev-vortex-free-gauge}
    u_{ij}=+1
    \qquad
    \text{for every bond oriented from }A\text{ to }B .
\end{equation}
Using the unit-cell convention in \cref{eq:honeycomb-neighbour-convention}, the
Hamiltonian becomes
\begin{equation}
\label{eq:kitaev-real-space-vortex-free}
    H_0
    =
    \ii\sum_{\bm r}
    \left(
    J_z c_{A,\bm r}c_{B,\bm r}
    +J_x c_{A,\bm r}c_{B,\bm r+\bm a_1}
    +J_y c_{A,\bm r}c_{B,\bm r+\bm a_2}
    \right).
\end{equation}
After Fourier transformation, the only gauge-invariant spectral data are encoded
in
\begin{equation}
\label{eq:kitaev-bloch-function}
    f(\bm k)
    =
    J_z+J_x\ee^{\ii \bm k\cdot\bm a_1}
       +J_y\ee^{\ii \bm k\cdot\bm a_2} .
\end{equation}
The Majorana quasiparticle dispersion is
\begin{equation}
\label{eq:kitaev-majorana-dispersion}
    \epsilon(\bm k)=2\abs{f(\bm k)} .
\end{equation}

\begin{remark}[Factor of two convention]
The factor of two in \cref{eq:kitaev-majorana-dispersion} follows from the
Hamiltonian normalization in \cref{eq:kitaev-honeycomb-spin-hamiltonian}.  Some
references absorb this factor into the definition of \(J_\alpha\) or into the
canonical Majorana matrix \(A\).  The physically invariant statement is the
location of zeros of \(f(\bm k)\) and the relative gap scale.
\end{remark}

The gap closes if and only if there exists a momentum \(\bm k\) such that
\begin{equation}
\label{eq:kitaev-gap-closing-condition-f}
    J_z+J_x\ee^{\ii \bm k\cdot\bm a_1}
       +J_y\ee^{\ii \bm k\cdot\bm a_2}=0 .
\end{equation}
Geometrically, this means that three complex numbers of lengths
\(J_x,J_y,J_z\) can form a triangle.  Therefore the gapless condition is
\begin{equation}
\label{eq:kitaev-triangle-inequalities}
    J_x\leq J_y+J_z,
    \qquad
    J_y\leq J_z+J_x,
    \qquad
    J_z\leq J_x+J_y .
\end{equation}
Inside this region, the spectrum has two Dirac cones.  Outside it, one coupling
dominates and the spectrum is gapped.

\begin{proposition}[Bulk phase criterion]
For positive couplings, the vortex-free Majorana spectrum is gapless precisely
inside the triangle region \cref{eq:kitaev-triangle-inequalities}.  It is gapped
in the three regions
\begin{equation}
\label{eq:kitaev-a-phase-regions}
    J_x>J_y+J_z,
    \qquad
    J_y>J_z+J_x,
    \qquad
    J_z>J_x+J_y .
\end{equation}
\end{proposition}

\begin{proof}
The only possible gap closing is \(f(\bm k)=0\).  Since the phases
\(\ee^{\ii \bm k\cdot\bm a_1}\) and
\(\ee^{\ii \bm k\cdot\bm a_2}\) can vary over the unit circle subject to the
Brillouin-zone periodicity, the equation \(f(\bm k)=0\) is equivalent to closing
a triangle with side lengths \(J_x,J_y,J_z\).  Such a triangle exists if and
only if the three triangle inequalities hold.  If one inequality is violated,
\(\abs{f(\bm k)}\) has a strictly positive lower bound.
\end{proof}

At the isotropic point
\begin{equation}
\label{eq:kitaev-isotropic-point}
    J_x=J_y=J_z,
\end{equation}
the system lies in the gapless phase.  The low-energy theory consists of two
cones of massless Majorana fermions coupled to a gapped \(\ZZ_2\) gauge field.
Since the gauge field is static in the exactly solvable model, the low-energy
fermions are free within each flux sector, but flux excitations remain massive.

\section{Gapped Abelian phases and the toric-code limit}
\label{sec:kitaev-abelian-phases-toric-code-limit}

The three gapped regions in \cref{eq:kitaev-a-phase-regions} are usually called
\(A_x\), \(A_y\), and \(A_z\).  In the \(A_z\) phase, for example,
\begin{equation}
\label{eq:kitaev-az-limit-condition}
    J_z\gg J_x,J_y .
\end{equation}
The unperturbed Hamiltonian consists of decoupled \(z\)-dimers.  Treating the
\(x\)- and \(y\)-bond interactions perturbatively, the first nontrivial
plaquette-dependent term appears at fourth order and has the form
\begin{equation}
\label{eq:kitaev-effective-toric-code}
    H_{\rm eff}^{(4)}
    =
    -J_{\rm eff}\sum_p W_p+\text{constant},
    \qquad
    J_{\rm eff}\sim \frac{J_x^2J_y^2}{J_z^3} .
\end{equation}
More precisely, the numerical prefactor depends on the perturbative convention,
but the structure is fixed: the low-energy Hamiltonian is a \(\ZZ_2\)
commuting-projector theory generated by plaquette fluxes.

This shows that the gapped \(A\) phases are Abelian topological phases in the
same universality class as the toric code.  The elementary excitations can be
identified as
\begin{enumerate}[label=(\roman*)]
    \item \(\ZZ_2\) fluxes, corresponding to plaquettes with \(W_p=-1\);
    \item fermionic matter excitations from the Majorana band;
    \item bound states of fluxes and Majorana fermions.
\end{enumerate}
Their fusion and braiding reduce, in the deep Abelian limit, to the toric-code
anyon structure.

\begin{remark}[Why the toric code reappears]
The toric code describes a deconfined \(\ZZ_2\) gauge theory with gapped charges
and fluxes.  The honeycomb model in a strongly anisotropic gapped phase has the
same long-distance topological order, but its microscopic degrees of freedom are
spins on sites rather than qubits on edges.  The emergence of the toric code is
therefore an infrared statement, not an exact microscopic equality at generic
couplings.
\end{remark}

\section{Time-reversal breaking and the non-Abelian phase}
\label{sec:kitaev-non-abelian-phase}

The gapless phase is protected by the combination of lattice symmetries and
time-reversal symmetry.  A small magnetic field
\begin{equation}
\label{eq:kitaev-magnetic-field-perturbation}
    V_h=-\sum_i\left(h_x\sigma_i^x+h_y\sigma_i^y+h_z\sigma_i^z\right)
\end{equation}
breaks time reversal.  The magnetic field itself destroys exact solvability, but
its leading nonvanishing effect within perturbation theory is a solvable
three-spin interaction of the schematic form
\begin{equation}
\label{eq:kitaev-effective-three-spin-term}
    H_\kappa
    =
    -\kappa\sum_{\langle i j k\rangle}
    \sigma_i^\alpha\sigma_j^\beta\sigma_k^\gamma,
    \qquad
    \kappa\propto \frac{h_xh_yh_z}{J^2},
\end{equation}
where the ordered triples \(\langle i j k\rangle\) are length-two paths on the
honeycomb lattice and \(\alpha,\beta,\gamma\) are fixed by the adjacent bond
types.  In the Majorana representation this term adds next-nearest-neighbour
Majorana hopping.  It opens a mass gap at the Dirac cones.

The resulting gapped phase has a nonzero Majorana Chern number,
\begin{equation}
\label{eq:kitaev-majorana-chern-number}
    \nu=\pm1,
\end{equation}
with the sign determined by the sign of \(\kappa\).  The vortices of the
\(\ZZ_2\) gauge field bind Majorana zero modes.  Consequently, vortices become
non-Abelian anyons of Ising type.

The Ising-anyon fusion rules are
\begin{equation}
\label{eq:ising-anyon-fusion-rules}
    \sigma\times\sigma=1+\psi,
    \qquad
    \sigma\times\psi=\sigma,
    \qquad
    \psi\times\psi=1 .
\end{equation}
Here \(\sigma\) denotes a vortex carrying a Majorana zero mode, and \(\psi\) is
the fermionic Majorana excitation.  The quantum dimensions are
\begin{equation}
\label{eq:ising-anyon-quantum-dimensions}
    d_1=1,
    \qquad
    d_\psi=1,
    \qquad
    d_\sigma=\sqrt{2} .
\end{equation}
The noninteger quantum dimension of \(\sigma\) is the hallmark of non-Abelian
statistics.

\begin{remark}[Relation to chiral topological superconductors]
After fixing a gauge sector, the itinerant Majoranas form a two-dimensional
class-D superconductor.  The time-reversal-breaking mass converts the gapless
Majorana semimetal into a chiral Majorana topological superconductor.  Coupling
this band topology to the dynamical \(\ZZ_2\) fluxes produces the non-Abelian
Ising topological order.
\end{remark}

\section{Fluxes, vortices, and physical projection}
\label{sec:kitaev-fluxes-physical-projection}

The Majorana solution contains two levels of description.  First, one fixes a
configuration of link variables \(u_{ij}\) and diagonalizes the free Majorana
Hamiltonian \(H[u]\).  Second, one projects the resulting states onto the
physical spin Hilbert space using \(P_{\rm phys}\).  This projection has two
important consequences.

\begin{enumerate}[label=(\roman*)]
    \item Gauge-equivalent link configurations describe the same physical flux
    sector.  A local sign flip generated by \(D_i\) changes adjacent \(u_{ij}\)
    and \(c_i\), but leaves spin observables invariant.
    \item On a finite torus, the physical projection fixes a relation between
    the gauge-field sector and the fermion parity of the matter Majoranas.  Thus
    not every free-Majorana eigenstate in the enlarged Hilbert space represents
    a physical spin state.
\end{enumerate}

For bulk thermodynamic properties, the projection usually does not change the
local dispersion or the phase boundaries.  For exact finite-size spectra and
for topological degeneracy on a torus, however, it is essential.

A useful way to summarize the structure is
\begin{equation}
\label{eq:kitaev-sector-label-summary}
    \text{physical sector}
    =
    \frac{\text{link variables }u_{ij}=\pm1}{\text{local }\ZZ_2\text{ gauge transformations}}
    +
    \text{projected Majorana matter states} .
\end{equation}
Gauge-invariant fluxes \(W_p\) and noncontractible Wilson loops label the gauge
part, while the free Majorana occupation numbers label the matter part.

\section{Spin correlations and fractionalization}
\label{sec:kitaev-spin-correlations-fractionalization}

A single spin operator is not diagonal in the flux sector.  From
\(\sigma_i^\alpha=\ii b_i^\alpha c_i\), acting with \(\sigma_i^\alpha\) changes
link variables adjacent to site \(i\).  Equivalently, in the spin language it
creates or moves a pair of fluxes on the two plaquettes adjacent to the
\(\alpha\)-bond touching site \(i\).  Therefore ordinary spin correlations are
strongly constrained.

In the exactly solvable zero-field model, the equal-time two-spin correlator
vanishes beyond nearest neighbours:
\begin{equation}
\label{eq:kitaev-spin-correlation-range}
    \langle \sigma_i^\alpha\sigma_j^\beta\rangle=0
    \qquad
    \text{unless } i,j \text{ are nearest neighbours on an }\alpha\text{-bond}
    \text{ and }\alpha=\beta .
\end{equation}
This short-ranged spin correlation is not a sign of a trivial phase.  It reflects
spin fractionalization: a physical spin flip fractionalizes into an itinerant
Majorana excitation and a pair of \(\ZZ_2\) fluxes.

By contrast, gauge-invariant Majorana bilinears and flux operators can encode
long-distance information about the Majorana band structure.  In the gapless
phase, the itinerant Majoranas have algebraic correlations, but local spin
operators couple to them only together with flux excitations.

\section{Comparison with the toric code and quantum double models}
\label{sec:kitaev-comparison-toric-code-quantum-double}

The following table summarizes the relation between the solvability mechanisms
of the previous two chapters and the honeycomb model.

\begin{center}
\begin{adjustbox}{max width=\textwidth}
\begin{tabular}{@{}L{0.24\textwidth}L{0.28\textwidth}L{0.32\textwidth}@{}}
\toprule
Model class & Exact-solvability mechanism & Excitations \\
\midrule
Toric code &
Commuting stabilizers \(A_s,B_p\) &
Static Abelian anyons \(1,e,m,\epsilon\) \\
\addlinespace
Quantum double \(D(G)\) &
Commuting gauge and flatness projectors \(A_v,B_p\) &
Electric, magnetic, and dyonic anyons labeled by \((C,\rho)\) \\
\addlinespace
Kitaev honeycomb model &
Static \(\ZZ_2\) gauge sectors plus free Majorana matter &
Fluxes, itinerant Majorana fermions, and in the chiral phase non-Abelian Ising anyons \\
\bottomrule
\end{tabular}
\end{adjustbox}
\end{center}

The toric code and quantum double models are fixed-point commuting-projector
Hamiltonians.  The Kitaev honeycomb model is instead a microscopic spin model
with dispersive fractionalized matter.  Its gapped Abelian regions flow to
toric-code topological order at long distances, while its time-reversal-broken
gapless region flows to Ising topological order.

\section{Summary}
\label{sec:kitaev-honeycomb-summary}

The Kitaev honeycomb model is exactly solvable because the spin Hamiltonian can
be rewritten as free Majorana fermions moving in a static \(\ZZ_2\) gauge field.
The essential steps are:
\begin{enumerate}[label=(\roman*)]
    \item enlarge each spin into four Majoranas,
    \begin{equation}
    \label{eq:kitaev-summary-majorana-rep}
        \sigma_i^\alpha=\ii b_i^\alpha c_i,
        \qquad
        D_i=b_i^xb_i^yb_i^zc_i=+1;
    \end{equation}
    \item define conserved link variables,
    \begin{equation}
    \label{eq:kitaev-summary-link-variable}
        u_{ij}=\ii b_i^\alpha b_j^\alpha,
        \qquad
        u_{ij}=\pm1;
    \end{equation}
    \item reduce the Hamiltonian in every gauge sector to
    \begin{equation}
    \label{eq:kitaev-summary-free-majorana}
        H[u]=\ii\sum_{\langle ij\rangle_\alpha}J_\alpha u_{ij}c_i c_j;
    \end{equation}
    \item identify gauge-invariant fluxes
    \begin{equation}
    \label{eq:kitaev-summary-flux}
        W_p=\prod_{\langle ij\rangle\in\partial p}u_{ij};
    \end{equation}
    \item diagonalize the vortex-free Majorana band with dispersion
    \begin{equation}
    \label{eq:kitaev-summary-dispersion}
        \epsilon(\bm k)=2\abs{J_z+J_x\ee^{\ii\bm k\cdot\bm a_1}
        +J_y\ee^{\ii\bm k\cdot\bm a_2}} .
    \end{equation}
\end{enumerate}
The bulk phase diagram follows from the triangle inequalities
\cref{eq:kitaev-triangle-inequalities}.  One obtains gapped Abelian phases when
one coupling dominates, a gapless Majorana spin liquid inside the triangle, and
a non-Abelian Ising-anyon phase when the gapless phase is gapped by a
small time-reversal-breaking perturbation.
